\documentclass[12pt,english]{article}
\usepackage[T1]{fontenc}
\usepackage[latin9]{inputenc}
\usepackage[toc,page]{appendix}
\usepackage{geometry}
\geometry{verbose,tmargin=1.5in,bmargin=1.5in,lmargin=1.5in,rmargin=1.5in}
\usepackage{bm}
\usepackage{amstext}
\usepackage{comment}
\usepackage{amsmath}
\usepackage{hyperref}
\hypersetup{
     colorlinks = true,
     citecolor = navy,
     linkcolor = navy,
     urlcolor=navy
}
\usepackage{adjustbox}
\usepackage[space]{grffile}
\usepackage{setspace}
\usepackage[normalem]{ulem}
\usepackage{nicematrix}
\usepackage{xcolor}
\usepackage{soul}
\usepackage{arydshln}
\usepackage{dsfont}
\usepackage{refstyle}
\usepackage[capitalize,nameinlink]{cleveref}
\usepackage{setspace}
\usepackage{bbm}
\usepackage[tableposition=t]{caption} % ensures caption of longtables is normalsized.
\usepackage{mathpazo} % add Palatino font
\usepackage{longtable}
\usepackage{amssymb}
\usepackage{amsfonts}
\usepackage{amsmath}
%\usepackage[margin=1in]{geometry} % to set margins
\usepackage{scalefnt}
\usepackage{wrapfig}
\usepackage{graphicx}
\usepackage{graphics}

%\usepackage{subfigure}
\usepackage{lscape}
\usepackage{pdflscape}
\usepackage{multirow}
\usepackage{array}
\usepackage{rotating}
\usepackage{bibunits}
\usepackage{fix-cm}
\usepackage{float}
\floatplacement{table}{H}
\usepackage{epsfig}
\usepackage{epstopdf}
\usepackage{refstyle}
\usepackage{tikz}
\usepackage{booktabs}
\usepackage{tabularx}
\usepackage{array}
\newcolumntype{x}[1]{>{\raggedleft\arraybackslash}p{#1}}%
\newcolumntype{d}[1]{>{\centering\arraybackslash}p{#1}}%
\newcolumntype{y}[1]{>{\raggedright\arraybackslash}p{#1}}%
\usepackage{verbatim}
%\epstopdfsetup{suffix=}
\setcounter{MaxMatrixCols}{10}
\usepackage[longnamesfirst,authoryear]{natbib}
\newenvironment{proof}[1][Proof]{\noindent \textbf{#1.}}{\rule{0.5em}{0.5em}}
\setlength{\topmargin}{-0.80in} \setlength{\textheight}{9.50in}
\setlength{\evensidemargin}{.0in} \setlength{\oddsidemargin}{.0in}
\setlength{\textwidth}{6.85in}
\setlength{\voffset}{0.4cm}
\setlength{\hoffset}{-0.5cm}
\usepackage{setspace}
\usepackage{afterpage}
\usepackage[flushleft]{threeparttable}
\usepackage{tabulary}
%\usepackage[titletoc,title]{appendix}
\usepackage[bottom]{footmisc}
\usepackage{color}
	\definecolor{maroon}{RGB}{128,0,0}
	\definecolor{navy}{RGB}{0,0,165}
	\definecolor{lb}{RGB}{150,165,190}
	\definecolor{bluenice}{rgb}{0.25,0.51,0.76}
	\definecolor{rednice}{rgb}{0.85,0.26,0}
	\definecolor{orangenice}{HTML}{EB811B}
	\definecolor{greennice}{HTML}{14B03D}
\usepackage{colortbl}
\usepackage{caption}
\usepackage{subcaption}
\usepackage{placeins}
\usepackage{nicefrac}
%\usepackage{chngcntr}
\newtheorem{assumption}{Assumption}
\newtheorem{definition}{Definition}
\usetikzlibrary{matrix,decorations.pathreplacing}

\newcommand{\vy}{\underline{y}}
\newcommand{\ve}{\underline{\epsilon}}
\newcommand{\A}[1]{\mathbold{A_{#1}}}
\newcommand{\D}[1]{\mathbold{D_{#1}}}
\newcommand{\C}[1]{\mathbold{C_{#1}}}
\newcommand{\B}[1]{\mathbold{B_{#1}}}
\newcommand{\R}[1]{\mathbold{R_{#1}}}

%\newtheorem{definition}{Definition}
\begin{document}
\title{Some possible checks to see what our influence correction is doing}
\author{N}
\maketitle


\paragraph{Purpose:} To come up with some checks to see what is going on with our influence calculations.

\paragraph{First check: Expenditure Shares vs. Influence Impacts}

We know that CPI is given by:

\begin{equation}
	\log(P)_{CPI,t} = \underline{\omega}_{t}'\log(\underline{P}_t)\label{eqn:CPI}
\end{equation}
where $\underline{x}$ is an $N$ x $1$ vector of variable $x$ (either $P$ or $\omega$) where $N$ is the number of products and $\underline{\omega}$ is a vector of expenditure shares in consumption (and sums to 1).

We also calculated influence using an $N$ by $N$ matrix I called $\mathbf{\Theta}$ which is calculated as:

\[
	\mathbf{\Theta}\equiv (I_N - diag(\underline{\gamma})\mathbf{\Omega})^{-1}diag(1-\underline{\gamma})
\]
where $\underline{\gamma}$ is the vector of intermediate shares of sales for each sector and $\mathbf{\Omega}$ is the I-O matrix.

This matrix was used to apportion value added price $\underline{P}_{VA,t}$ changes to output prices $\underline{P}_t$. When we decided some value added price changes were due to demand, we used this matrix to translate that to changes in output prices.

Here is what I would like done:

\begin{enumerate}
	\item Calculate $\underline{\omega}_{VA}\equiv \mathbf{\Theta}'\underline{\omega}$ (don't forget the transpose). This vector translates product value added ``prices'' to CPI the same way Equation \ref{eqn:CPI} translates product output prices to CPI.
	\item Pick 3 time periods. One early in the sample, one near 1981 and one near the end of the sample. Then for each do the following
	\item Order the sectors by highest to lowest expenditure shares in CPI: $\underline{\omega}_{t}$. Apply the \underline{same} ordering to the value added equivalent $\underline{\omega}_{VA}$. That is if both are in the same dataframe, we are ordering the rows the same way by the sectors/products with the largest $\omega_{i,t}$ to the smallest.
	\item Plot both series with products/sectors on the horizontal and the value of $\omega_{i,t}$ and $\omega_{i,VA,t}$ on the vertical. Each concept has its own lines.
\end{enumerate}

\emph{What we are looking for:} Based on Ben's example, we are wanting to see if these are very different. 

{\bf Note:} If it is hard to see the differences between the two lines, take the log of $\omega_{i,t}$ and $\omega_{i,VA,t}$ and replot. Both should be positive so this should be able to be done. 

\paragraph{Second check: Distribution of Demand Shocks and Supply Shocks}

Next we should see if we are in a situation where the demand shocks and supply shock \emph{values} look similar. To do this, for the same 3 time periods as above, I want you to take the variables you calculated:

\begin{align*}
	\log(P_{VA,t})^{D=1} &\equiv \log(\underline{P}_{VA,t}) \circ \tilde{\underline{D}}_t\\
	\log(P_{VA,t})^{D=0} &\equiv \log(\underline{P}_{VA,t})\circ (1-\tilde{\underline{D}}_t) 
\end{align*}
where $\underline{P}_{VA,t}$ is a vector containing the value added ``price'' of value added in each sector at time $t$ and $\tilde{\underline{D}}_t$ is a vector indicating whether the sector has faced a demand shock or a supply shock (it is 1 or 0).

Here is what I suggest we do:

\begin{enumerate}
	\item Remove any sectors in $\log(P_{VA,t})^{D=1}$ where $\tilde{\underline{D}}=0$. Before we kept these and made them 0. Now I want them removed.
	\item Do the same for $\log(P_{VA,t})^{D=0}$ when $\tilde{\underline{D}}_t=1$
	\item Do a histogram or kernal density plot of both of these. Ideally it would be transparent so we can see both of the distributions like the figure below:
	
	\includegraphics[width=8cm]{graph_to_use.png}
	
	where we underline{don't} need the vertical lines indicating the mean. All I want is the transparent shading if possible.
	\item One other thing: if doing a kernal density be sure to account for the fact that we want to see if we have \emph{more} of supply shocked sectors or demand shocked sectors. If the kernal density function makes the area under the density sum to one, then we don't want to use that function -- a histogram would be better
\end{enumerate}

\emph{What we are looking for:} Is it the case that supply and demand densities are the same or not? If they are the same, then unless we have super skewed expenditure/influence weights, we will see very similar supply/demand contributions regardless of the weights used.

\paragraph{Third Check: Top 10 contributors}

Finally, I propose we plot the contribution of the 10 sectors/products with the largest contribution (or do a table of them).

We calculate contribution as follows:

\[
	contrib_t =\frac{1}{P_{CPI,t}}\underline{\omega}_{VA,t}'diag(P_{VA,t})
\]
 
 $contrib_t$ should be an $N$ x $1$ vector \emph{and} should sum to 1. {\bf Let me know if it does not.}
 
 Once it sums to 1, we can plot/table the top 10 values of it. If we are plotting these, then let's do a bar graph with 1 color for mostly supply sectors and another for mostly demand. If we are doing a table, then include the mostly demand/mostly supply column please!

\emph{What we are looking for:} This calculates the interaction of the highest weights (influence) with the size of the shock. Basically we are seeing if anything pops out here. Are the top 10 sectors contributing a large amount to CPI (e.g. 20\% or more) or are the contributions very diffuse?


\paragraph{Additional thing to do on graph 2 if time}

Let's replicate the second graph with Shapiro shocks. The way to do this is to replace the value added weights with expenditure weights and the value added prices with output prices. 


\begin{align*}
	\log(P_{t})^{D^{Shapiro}=1} &\equiv \log(\underline{P}_{t}) \circ \tilde{\underline{D}}_t^{Shapiro}\\
	\log(P_{t})^{D^{Shapiro}=0} &\equiv \log(\underline{P}_{t})\circ (1-\tilde{\underline{D}}_t^{Shapiro}) 
\end{align*}

And repeat the steps with these objects.

Thanks!

\end{document}